% ===================================================================
%  Section II: Background and Preliminaries
% ===================================================================

This section reviews the technical foundations used throughout the paper.

\subsection{Markov Decision Process and DQN}

We consider a discounted Markov decision process (MDP)
$(\states, \actions, P, r, \gamma)$ with state space $\states$, finite action
space $\actions$, transition kernel $P$, bounded reward function $r$, and
discount factor $\gamma \in [0,1)$. A policy $\pi : \states \to \actions$
selects actions, and the action-value function under $\pi$ is
$Q^\pi(s,a) = \mathbb{E}_\pi\bigl[\sum_{t=0}^{\infty} \gamma^t r_t
\mid s_0 = s, a_0 = a\bigr]$.

Deep Q-Network (DQN)~\cite{mnih2015dqn} approximates $Q^*$ with a neural
network $Q_\theta$ and uses experience replay and a periodically updated target
network $Q_{\bar\theta}$ to stabilize training. The loss for a sampled
transition $(s, a, r, s', d)$ is
\begin{equation}\label{eq:dqn-loss}
  \ell(\theta) = L\bigl(Q_\theta(s,a) - y\bigr),
\end{equation}
where $d \in \{0,1\}$ is the terminal flag and the Bellman target is
\begin{equation}\label{eq:bellman-target}
  y = r + \gamma (1 - d) \max_{a' \in \actions} Q_{\bar\theta}(s', a').
\end{equation}
Double DQN~\cite{hasselt2016double} decouples action selection from target
evaluation:
\begin{equation}\label{eq:ddqn-target}
  a^* = \arg\max_{a'} Q_\theta(s', a'), \quad
  y = r + \gamma (1 - d)\, Q_{\bar\theta}(s', a^*).
\end{equation}

Throughout this paper, the loss function $L$ is the SmoothL1 (Huber)
loss~\cite{huber1964robust}:
\begin{equation}\label{eq:smoothl1}
  L(\delta) =
  \begin{cases}
    \tfrac{1}{2}\delta^2 & \text{if } |\delta| < 1,\\
    |\delta| - \tfrac{1}{2} & \text{otherwise.}
  \end{cases}
\end{equation}

\subsection{Offline Reinforcement Learning}

Offline RL~\cite{levine2020offline, fujimoto2019offpolicy} trains policies from
a fixed dataset $\dataset = \{(s_i, a_i, r_i, s'_i, d_i)\}_{i=1}^N$ without
additional environment interaction. This eliminates exploration but introduces
distributional challenges: value overestimation for out-of-distribution
actions~\cite{kumar2020cql} and sensitivity to dataset quality. From a
verification perspective, the fixed dataset creates a clear commitment point:
the verifier can require that all training computations reference transitions
from a committed dataset.

\subsection{Zero-Knowledge Proofs}\label{subsec:zkp}

A zero-knowledge proof system~\cite{thalerbook} allows a prover~$\mathcal{P}$
to convince a verifier~$\mathcal{V}$ that a statement~$x$ is true with respect
to a witness~$w$ satisfying a relation~$R(x,w) = 1$, without revealing~$w$
beyond the validity of the statement. A proof system is:
\begin{itemize}[leftmargin=1.5em]
  \item \emph{Complete}: an honest prover with a valid witness always
  convinces the verifier.
  \item \emph{Sound}: no computationally bounded prover can convince the
  verifier of a false statement (except with negligible probability).
  \item \emph{Zero-knowledge}: the verifier learns nothing beyond the truth
  of the statement.
\end{itemize}

\noindent A \emph{succinct} proof system (zkSNARK/zkSTARK) additionally
guarantees that verification time and proof size are sublinear (or polylogarithmic)
in the computation size~\cite{groth2016zksnark, bensasson2014scalable}.

\subsection{Merkle Tree Commitments}\label{subsec:merkle}

A Merkle tree~\cite{merkle1987} is a binary hash tree in which each leaf
contains the hash of a data element and each internal node contains the hash
of its two children. The value $\mroot$ serves as a compact commitment to all
leaves. To prove membership of leaf $\ell_i$, the prover provides a sibling
path $(h_1, h_2, \ldots, h_d)$ with direction bits, and the verifier
recomputes the root from the leaf through $d$ hash steps. Membership
verification requires $O(d) = O(\log N)$ hash evaluations for $N$ leaves.

In our system, each replay transition $\tau_i = (s_i, a_i, r_i, s'_i, d_i)$
is serialized into a canonical fixed-point encoding, hashed with SHA-256, and
inserted as a leaf. The dataset commitment additionally binds the Merkle root
to a manifest hash, audit-report hash, raw trajectory hash, and collection-log
hash (when available), creating a \emph{provenance-bound} commitment.

\subsection{SP1 Zero-Knowledge Virtual Machine}\label{subsec:sp1}

SP1~\cite{succinct2024sp1} is a RISC-V-based general-purpose zkVM that allows
developers to write provable programs in Rust. SP1 uses a STARK-based proof
system with recursive composition: the program executes inside the VM, and the
execution trace is proved segment by segment. For on-chain verification, proofs
can be compressed to a Groth16 SNARK~\cite{groth2016zksnark}.

We chose SP1 as the proof backend because our relations combine SHA-256
hashing, Merkle path traversal, fixed-point integer arithmetic, branching
logic, and schema validation---a workload well suited to a general-purpose
zkVM. Each relation is structured as a \emph{host/guest/shared} Rust workspace:
\begin{itemize}[leftmargin=1.5em]
  \item \textbf{Shared crate}: typed public inputs, private witnesses,
  fixed-point helpers, and relation-check logic.
  \item \textbf{Guest crate}: reads inputs inside the zkVM and calls the
  shared relation verifier.
  \item \textbf{Host crate}: loads JSON fixtures, selects execute or prove
  mode, verifies proofs, and records compact metrics.
\end{itemize}

\subsection{Fixed-Point Arithmetic}\label{subsec:fixedpoint}

All arithmetic in the relation layer uses deterministic fixed-point integer
encoding with a global scale factor $\fpscale = 1000$. A real value $v$ is
represented as $\lfloor v \cdot \fpscale \rceil \in \mathbb{Z}$, and all
additions, multiplications (with rescaling), comparisons, and hash inputs
operate on these integer values. This eliminates floating-point
nondeterminism between Python semantic oracles and Rust/SP1 guests, ensuring
that the same inputs produce identical outputs across both verification
paths.
